% PNP RESEARCH REPORT::RECOVERY-STRENGTHEN2::2026-10-08T20:47:00Z
\documentclass[11pt]{article}
\usepackage[T1]{fontenc}
\usepackage[utf8]{inputenc}
\usepackage[margin=1in]{geometry}
\usepackage{amsmath,amssymb,amsthm,booktabs,hyperref}
\usepackage{xurl}
\emergencystretch=3em
\newtheorem{proposition}{Proposition}
\title{RECOVERY-STRENGTHEN2: Returning Phases and the Exact Access Budget}
\author{P versus NP Research Project}
\date{8 October 2026, 20:47 UTC}
\begin{document}\maketitle

\section*{Status and substantive recovery}
P versus NP remains open. The original reports already prove conditional
elimination of the polynomial SAT exponent $K$, allow typed cycles whose
complete cycle contracts, and allow shared acquisition and charged implicit
input access. Requiring every internal edge to shrink, every child to reacquire
shared setup, or every logical same-scale input to occupy a fresh linear buffer
would discard legitimate candidates. None of those requirements is necessary.
The missing implementation is still not supplied by the inspected sources.

This run independently verifies the previously unpublished
RECOVERY-STRENGTHEN1 bounded-phase recovery and derives an exact two-type
near-linear feasibility test. This is reusable accounting, not a newly
constructed all-input SAT solver. In particular, it identifies precisely when
access overhead consumes the contraction gain in this envelope.

\section*{Corpus, provenance, and reconciliation}
Read current root \nolinkurl{main.tex}, \nolinkurl{main4.tex},
\nolinkurl{main5.tex}, \nolinkurl{AGENTS.md}, \nolinkurl{CODEX.md}, the
repository research skill and protocol, FORMAL-SPINE1, and the latest
CONSTRUCTION-STRENGTHEN2 aggregate-volume report. The four original source
reports, read from GitHub rather than inferred from later summaries, are:
\begin{description}
\item[RECURCLOSE1.] \nolinkurl{research/agents/lspv-2/2026-10-04_1604_LSPV-2_RECURCLOSE1.tex},
blob \nolinkurl{5c77ed64e79493883a52d0b5425dd195ba17b3cf}.
\item[CYCLEGAIN1.] \nolinkurl{research/agents/lspv-3/2026-10-04_1937_LSPV-3_CYCLEGAIN1.tex},
blob \nolinkurl{aa3859c364c0c6eb15a653f0218105cf9eb1917a}.
\item[STREAMCOMP1.] \nolinkurl{research/agents/lspv-3/2026-10-04_2135_LSPV-3_STREAMCOMP1.tex},
blob \nolinkurl{adc67bbc9f5d0255defc770cd48117dece0741d7}.
\item[ORBITBATCH1.] \nolinkurl{research/agents/lspv-3/2026-10-05_0235_LSPV-3_ORBITBATCH1.tex},
blob \nolinkurl{766891e34f5151ed542e407d0f8bfee63322c6ba}.
\end{description}
RECOVERY-STRENGTHEN1 and its contemporaneous log are recovered unchanged
from an accessible earlier workspace; this run publishes them with a new
independent review. Their historical ``not written'' and ``review pending''
fields are preserved as historical facts, superseded by this run's handoff.
The original four reports remain their source of authority.

Adopt RECURCLOSE1's recurrence and conditional $K$ elimination; CYCLEGAIN1's
typed contraction; STREAMCOMP1's explicit access-cost obligations; and
ORBITBATCH1's shared setup and address deduplication. Reject use of the old
unrestricted $t+s<2\cos(\pi/7)$ shorthand in RECURCLOSE1. Preserve the
current root's OPEN status and ownership; reconcile only the authorized
\nolinkurl{main4.tex} and \nolinkurl{main5.tex} additions. This is a
clarification and specialization of earlier work, not a claim to have
discovered conditional recursion closure for the first time.

\section*{Exact target, quantifiers, and realization}
The minimum useful result is an auditable contract for a whole returning
phase, plus a test of whether its proposed costs fit the near-linear target.
Fix one deterministic uniform resumable program $W$, constants $0<h<1$
and $0\le p<1$, and a fixed finite base threshold. For every encoded SAT
input $x$ of bit length $N$ in the claimed domain, and for every reached
recursive entry above that threshold, require the bounds below with
uniform $o(1)$ terms. The domain must contain all recursively generated
instances; the endpoint application requires all SAT inputs.

Use the random-access Turing-machine model matched to Fortnow's endpoint:
read-only entry input, separately charged work storage, and no unbounded
unit-cost arithmetic. The SAME\_REALIZATION object is the full execution
of this one $W$, including its controller, evaluations, child generation,
classification, failed attempts, repeat executions, replay, indexing,
buffers and suspended ancestor frames. No independent time witness and
space witness may be combined. No free semantic oracle is part of $W$.

\section*{Recovered returning-phase lemma}
A phase is the work between entry to a distinguished recursive type and
return, with its calls back to the distinguished type treated as child
calls. It may contain a fixed bounded chain of same-scale typed routines.
The bound on the chain alone is insufficient: all its executed work must
fit the following aggregate premises. Let $m_j\ge1$ be lengths of actual
executed recursive children, counted with multiplicity, and assume
\[
 L(N)\le N^{1+o(1)},\qquad
 V(N):=\sum_j m_j\le N^{1+o(1)},\qquad
 M(N):=\max_jm_j\le N^{h+o(1)}.
\]
The maximum of an empty child multiset is zero. All non-child phase
workspace, including peak scratch, all live typed frames and access
evaluators, is at most $N^{p+o(1)}$. At most one explicit child input of
size $N^{h+o(1)}$ is retained at a time. Other retained children, address
tables, descriptions, caches and regeneration state must fit the stated
phase workspace. All recursive entries obey the same contract and answers
are combined exactly.

\begin{proposition}[Conditional phase closure]
Under these premises the single depth-first realization satisfies
\[
 T(N)\le N^{1+o(1)},\qquad
 S(N)\le N^{\max(p,h)+o(1)}.
\]
The conclusion is independent of the exponent of a fixed polynomial SAT
decider used only in the finite base region.
\end{proposition}
\begin{proof}
Flatten the charged internal phase into one local operation. For fixed
$\epsilon>0$, uniformity gives a sufficiently large threshold and arbitrarily
small fixed $\delta>0$ with $L,V\le N^{1+\delta}$ and
$M\le N^{h+\delta}$. For strong induction with $T(m)\le C m^{1+\epsilon}$,
\[
 \sum_jT(m_j)\le C M^\epsilon V
 \le C N^{1+h\epsilon+\delta(1+\epsilon)}.
\]
Choose $\delta<\epsilon$ with
$h\epsilon+\delta(1+\epsilon)<\epsilon$ and $h+\delta<1$.
The resulting exponent slack absorbs both the child sum and $L(N)$
in $CN^{1+\epsilon}$ for all sufficiently large $N$. The finite base
region is absorbed in $C$. This proves the bound for the same algorithm
for every fixed $\epsilon$, rather than a sequence of different algorithms.
Along one active recursion path the sizes decrease by a fixed fractional
power, so its depth is $O(\log\log N)$. Summing all charged phase frames
and explicit child inputs yields the asserted space bound. All operations
within each phase already belong to its local and peak resource bounds.
\end{proof}

The older count contract $q\le N^{1-h+o(1)}$ implies the displayed volume
contract together with $M\le N^{h+o(1)}$. The volume form adopts the latest
construction report's stronger accounting. It does not license deduplicating
equal children that are in fact executed repeatedly. The independent review
verifies the older count-form flattening and two-type algebra; the volume
extension here is an explicit application of CONSTRUCTION-STRENGTHEN2.

\section*{Two-type test: when a same-scale phase helps}
Suppose the recovered typed analysis gives the sufficient exponent envelope
\[
 C_P\ge r_P,\quad C_P\ge u+C_Q,\quad
 C_Q\ge r_Q,\quad C_Q\ge v+hC_P,
 \qquad u,v\ge0,\quad 0<h<1.
\]
The $P\to Q$ edge is same-scale and the $Q\to P$ edge shrinks. The
exponents $u,v$ include executed multiplicity and access overhead, not
merely distinct labels. The local floors $r_P,r_Q$ include all local
work and any unshared acquisition. This is an algebraic envelope for a
source-proved recurrence, not a way to obtain such a recurrence from $P=NP$.

\begin{proposition}[Exact least feasible envelope]
For nonnegative floors and the displayed parameters, the componentwise
least feasible pair is
\[
 C_P^*=\max\left(r_P,u+r_Q,\frac{u+v}{1-h}\right),\qquad
 C_Q^*=\max(r_Q,v+hC_P^*).
\]
Consequently $C_P^*\le1$ exactly when
\[
 \boxed{r_P\le1,\qquad u+r_Q\le1,\qquad u+v\le1-h.}
\]
\end{proposition}
\begin{proof}
Every feasible pair satisfies $C_P\ge r_P$, $C_P\ge u+r_Q$, and
$C_P\ge u+v+hC_P$. These imply $C_P\ge C_P^*$ and then
$C_Q\ge C_Q^*$. Conversely the defining two last lower bounds on
$C_P^*$ give $u+C_Q^*\le C_P^*$. Thus the proposed pair is feasible.
The three tests for $C_P^*\le1$ follow directly.
\end{proof}

This is a PROVED/EXACT statement about inequalities. Its interpretation
as an algorithmic upper bound remains CONDITIONAL on source-visible
recurrences, termination and uniformity. It is not a lower bound on the
best SAT algorithm and not an equality claim for a particular implementation.
For a same-scale $Q$ routine with local floor $r_Q=1$, the near-linear
budget forces $u=0$: polynomially many independent linear-cost $Q$ calls
do not fit merely because each eventually returns smaller instances.
If $r_P=r_Q=1$, $u=0$, and $v=1-h+g$ for an additional charged access
exponent $g\ge0$, then
\[
 C_P^*=1+\frac{g}{1-h}.
\]
Any positive fixed $g$ therefore loses this particular near-linear envelope.
This does not rule out a tighter amortized recurrence or a different endpoint.
The contract is sensitive to how access is charged, not to a name such as
``streaming'' or ``double wrapper.''

\section*{Recovered premises versus missing implementations}
RECURCLOSE1 removes $K$ after a fixed shrinker is supplied; it does not
prove that every polynomial SAT algorithm yields that shrinker.
CYCLEGAIN1 permits a same-scale edge on a contracting cycle; it does not
make a positive-cost noncontracting cycle terminate cheaply. STREAMCOMP1
permits generated input bits, but all nested requests to an implicit parent
must be realized and charged. A fresh derived $\Theta(N)$-bit buffer
breaks sublinear local workspace. Aliasing an already supplied read-only
input causes no such fresh allocation; nonroot input storage is still
charged at its caller.

ORBITBATCH1 permits one shared acquisition, then address deduplication
and batched consumers. It does not remove actual output consumption.
Its field-domain normalization $N_{\rm field}=p^r$ is not the original
SAT input length. A bound measured in $M=N^\lambda$ with child exponent
$b$ gives root-scale child exponent $\lambda b$, not $b$. In particular
$\lambda=2,b=3/4$ is expansive on the root scale. Setup is a maximum
term in a correct charged recurrence when actually shared; repeatedly
invoking a setup implementation still incurs repeated cost.

\section*{Primary literature and exact endpoint}
Fortnow's \emph{Time-Space Tradeoffs for Satisfiability}, Theorem 3.1 and
Corollary 3.3, exclude one deterministic SAT machine using simultaneous
near-linear time and $N^{1-\varepsilon}$ space for fixed $\varepsilon>0$;
Section 3 establishes applicability to random-access machines. The input
measure is encoded formula length. Use the paper's constructibility
conditions and its actual model, not an arbitrary word-RAM with free
large arithmetic. For fixed $\rho=\max(p,h)<1$, the phase conclusion is
bounded by $N^{1-\varepsilon}$ space for some fixed positive $\varepsilon$.
Thus a universal implementation of the contract is excluded already;
the missing project bridge is deriving such an implementation from
$\mathrm{SAT}\in P$.
Source: \url{https://lance.fortnow.com/papers/files/npvsnl.pdf}.

Williams's ECCC TR23-082, Theorem 1.1, transfers an algorithm for CircuitSAT
on $n$-variable circuits of size $2^{\alpha n+o(n)}$, taking
$2^{\beta n+o(n)}$ time, to time $2^{(\beta-\alpha)n+o(n)}$ on
subexponential-size circuits, for $0<\alpha\le\beta$. Its stated model
requires nearly size-preserving time-to-circuit simulation; Section 3.2
addresses random access separately. This is a time theorem at specified
circuit sizes; it supplies no simultaneous sublinear workspace compiler
from arbitrary polynomial-time SAT. It does not close the present bridge.
Source: \url{https://eccc.weizmann.ac.il/report/2023/082/download}.

As an optimization analogy, discounted Bellman equations separate a
contractive value recursion from acquisition of its transition model.
Here substitution gives the scalar contraction
$F(z)=\max(r_P,u+r_Q,u+v+hz)$, whose Lipschitz constant is at most $h$.
The elementary fixed-point argument transfers exactly at this algebraic
level. It provides no SAT transition generator, semantic reduction or
workspace realization. Classification: ANALOGY / CANDIDATE for construction,
not a theorem transferring a control-system algorithm to SAT.
Primary teaching source: MIT 6.231, Bertsekas, lecture 15,
\url{https://ocw.mit.edu/courses/6-231-dynamic-programming-and-stochastic-control-fall-2015/resources/mit6_231f15_lec15/}.

\section*{Falsification, tests, and limitations}
The accompanying deterministic exact-rational test exhausts 1,200 parameter
tuples: $h\in\{1/4,1/2,3/4\}$, each local floor in
$\{0,1/2,1,3/2\}$, and each edge cost in $\{0,1/4,1/2,3/4,1\}$.
It verifies the inequalities, fixed-point equalities, the three near-linear
criteria, 4,800 strict coordinate-decrease controls, and 24,000 monotone
iteration steps. It separately checks nine access-budget cases, repeated
execution versus unique-label volume, and the source-normalization control.
The code records its own SHA-256 and enforces a 25-second alarm. The
symbolic proofs, not these finite checks, establish the algebra and
conditional closure. No SAT benchmark coverage, physical memory measurement
or measured bit-machine complexity is claimed.

DEAD/COPY/DELAY and RECOMP controls are addressed by counting actual work:
eight executions of one 16-bit child have volume 128, not 16. A positive
cost cycle with $h=1$ would demand $C\ge C+u+v$ and is rejected when
$u+v>0$. Even a zero-cost unit-gain envelope alone proves neither termination
nor elimination of $K$. Horn, dual-Horn, 2-SAT, affine and bounded-treewidth
inputs create no exception to these bookkeeping identities; no hardness
property is being attributed to those easy families. Renaming, padding and
auxiliary encodings require remeasuring the actual input bits and work.

The recurrence proof also works with oracle calls if their costs and
interfaces satisfy its premises. It therefore supplies no nonrelativizing
SAT bridge. Giving semantic or low-degree-extension answers for free
removes the very acquisition obligation being audited; source access is
not itself a barrier-avoidance proof. No Natural-Proofs circuit property
or cryptographic assumption is used. These controls reject resource
shortcuts; they do not establish an unrestricted SAT lower bound.

\section*{Handoff and remaining arrow}
The recovered phase lemma and algebra have independent review in
\nolinkurl{2026-10-08_RECOVERY-STRENGTHEN_phase_review.tex}.
The stronger aggregate phase statement is proved above from the current
construction report and also passed the independent publication review
addendum. Conditional $K$ elimination is preserved. The open
arrow is still
\[
 \mathrm{SAT}\in P\ \Longrightarrow\
 \text{one exact all-input, jointly favorable returning-phase implementation}.
\]
Next target: select one concrete STREAMCOMP/ORBITBATCH evaluator and prove
its nested bit-access trace in original SAT bit length, including the
whole-phase peak workspace and the three displayed time-budget inequalities.
Another abstract recurrence without that trace will not close this arrow.

Artifacts share prefix
\nolinkurl{research/agents/proof/2026-10-08_2047_RECOVERY-STRENGTHEN2}:
this report, \nolinkurl{-test.py}, \nolinkurl{-results.json}, and
\nolinkurl{-handoff.json}. Manuscript additions are under
\nolinkurl{research/manuscripts/RECOVERY-STRENGTHEN2/}.
% END_OF_VERIFIED_REPORT::RECOVERY-STRENGTHEN::RECOVERY-STRENGTHEN2
\end{document}
